Pour chacune des affirmations suivantes, indiquer si elle est vraie ou fausse. Justifier chaque réponse. Une réponse non justifiée ne rapporte aucun point.

	\begin{enumerate}[itemsep=1cm]
		\item L'espace est muni d'un repère orthonormé \Oijk{}.

		On considère :

		\begin{itemize}[label={\textbullet~}]
			\item la droite $(d)$ dont une représentation paramétrique est :
		\end{itemize}
		\[\begin{cases}
				x=t \\
			y=-1,5-t \\
			z=2-2 t
		\end{cases}\text { où } t \in \R; \]
		\begin{itemize}[label={\textbullet~}]
			\item les points A$(3 ; 0 ; 2)$ et B$(2 ; 1 ;-3)$;

			\item le plan $(\mathcal{P})$ d'équation : \quad $-x+y-5 z-0,5=0$.
		\end{itemize}

		\begin{enumerate}[itemsep=5mm]
			\item \textbf{Affirmation 1 :}

			Le plan $(\mathcal{P})$ est orthogonal à la droite (AB), et passe par le milieu du segment [AB].

			\item \textbf{Affirmation 2 :}

			Les droites $(d)$ et (AB) sont sécantes.

			\item On considère le point C de coordonnées $(1,5 ; -3 ; -1)$.

			\textbf{Affirmation 3 :}

			La mesure de l'angle $\widehat{\mathrm{ACB}}$, arrondie à $10^{-1}$, est égale à $70,5$\,\textdegree.
		\end{enumerate}

		\item Titouan et Clotilde participent à un escape game. Ils se retrouvent dans une salle possédant deux portes de sortie, A et B, chacune protégée par un digicode équipé de 8 touches portant des symboles différents :

		\begin{itemize}[label={\textbullet~}]
			\item le digicode de la porte A utilise un code de 3 symboles différents devant être saisis dans l'ordre;
			\item celui de la porte B utilise un code de 4 symboles différents qui peuvent être saisis dans n'importe quel ordre.
		\end{itemize}

		N'étant pas parvenus à obtenir d'indices, ils décident de procéder au hasard à la saisie des codes.

		Clotilde choisit un code pour la porte A, et Titouan choisit un code pour la porte B.

		On considère que les choix des codes sont équiprobables.

		\textbf{Affirmation 4 :}

		Titouan a plus de chances d'ouvrir sa porte que Clotilde.
	\end{enumerate}
